\chapter{Tabela de Intervalos de Pontuação do SARESP} \label{ap:apendice1}

O \textit{Intervalo de Pontuação do SARESP} é a forma da qual é contabilizado o nível de proficiência do aluno, dada a sua pontuação no exame. Os dados para a tabela estão disponível no site do SARESP no portal dos \href{https://dados.educacao.sp.gov.br/story/saresp}{Dados Abertos da Educação.} Em negrito, foram destacadas as notas do 9° ano do Ensino Fundamental (9°EF) e do 3° ano do Ensino Médio (3°EM), que foram as analisadas empiricamente no trabalho.

\begin{table}[hbt!]
    \centering
    \caption{Intervalos de Pontuação}

    \renewcommand{\arraystretch}{1.4}

    \textbf{Língua Portuguesa}

    \begin{tabular}{lccccc}
        \hline
         & 3ºEF & 5ºEF & 7ºEF & \textbf{9ºEF} & \textbf{3ºEM} \\
        \hline
        Abaixo do Básico & <125 & <150 & <175 & \textbf{<200} & \textbf{<250} \\
        Básico           & 125 a 174 & 150 a 199 & 175 a 224 & \textbf{200 a 274} & \textbf{250 a 299} \\
        Adequado         & 175 a 254 & 200 a 249 & 225 a 274 & \textbf{275 a 324} & \textbf{300 a 374} \\
        Avançado         & $\geq$225 & $\geq$250 & $\geq$275 & \textbf{$\geq$325} & \textbf{$\geq$375} \\
        \hline
    \end{tabular}

    \vspace{0.8cm}

    \textbf{Matemática}

    \begin{tabular}{lccccc}
        \hline
         & 3ºEF & 5ºEF & 7ºEF & \textbf{9ºEF} & \textbf{3ºEM} \\
        \hline
        Abaixo do Básico & <150 & <175 & <200 & \textbf{<225} & \textbf{<275} \\
        Básico           & 150 a 199 & 175 a 224 & 200 a 249 & \textbf{225 a 299} & \textbf{275 a 349} \\
        Adequado         & 200 a 249 & 225 a 274 & 250 a 299 & \textbf{300 a 349} & \textbf{350 a 399} \\
        Avançado         & $\geq$250 & $\geq$275 & $\geq$300 & \textbf{$\geq$350} & \textbf{$\geq$400} \\
        \hline
    \end{tabular}

    \caption*{Fonte: \href{https://dados.educacao.sp.gov.br/story/saresp}{Dados Abertos da Educação}, elaboração do autor.}
    \label{tab:intervalos_saresp}
\end{table}